
\documentclass[10pt]{article}
\usepackage[T1]{fontenc}
\usepackage[utf8]{inputenc}
\usepackage{lmodern}
\usepackage{amsmath,amssymb,amsthm,mathtools}
\usepackage{geometry}
\usepackage{enumitem}
\usepackage{tabularx}
\geometry{margin=0.86in}
\setlength{\parindent}{1.35em}
\setlength{\parskip}{0.18em}
\emergencystretch=3em
\setlist{leftmargin=2.2em,itemsep=0.12em,topsep=0.25em}
\newcommand{\A}{\mathbb A}
\newcommand{\Af}{\mathbb A^f}
\newcommand{\C}{\mathbb C}
\newcommand{\Gm}{\mathbb G_m}
\newcommand{\Q}{\mathbb Q}
\newcommand{\R}{\mathbb R}
\newcommand{\Z}{\mathbb Z}
\newcommand{\bS}{\mathbb S}
\newcommand{\GL}{\operatorname{GL}}
\newcommand{\Gp}{\operatorname{Gp}}
\newcommand{\Hom}{\operatorname{Hom}}
\newcommand{\Isom}{\operatorname{Isom}}
\newcommand{\Lie}{\operatorname{Lie}}
\newcommand{\Gal}{\operatorname{Gal}}
\newcommand{\Ker}{\operatorname{Ker}}
\newcommand{\ad}{\operatorname{ad}}
\newcommand{\dfn}{\operatorname{dfn}}
\newcommand{\id}{\operatorname{Id}}
\newcommand{\limproj}{\varprojlim}
\newcommand{\limind}{\varinjlim}
\newcommand{\MC}{M_{\C}}
\newcommand{\KMC}{K M_{\C}}
\newcommand{\whZ}{\widehat{\Z}}
\newcommand{\pii}{\pi_0}
\newcommand{\HasseH}{H^1}

\newcommand{\Spec}{\operatorname{Spec}}
\newcommand{\Res}{\operatorname{Res}}
\newcommand{\Norm}{\operatorname{N}}
\newcommand{\ab}{\mathrm{ab}}

\newcommand{\End}{\operatorname{End}}
\newcommand{\Tr}{\operatorname{Tr}}
\newcommand{\NS}{\operatorname{NS}}
\newcommand{\Aut}{\operatorname{Aut}}
\newcommand{\Id}{\operatorname{Id}}


\newcommand{\Sp}{\operatorname{Sp}}
\newcommand{\codim}{\operatorname{codim}}

\begin{document}


\begin{center}
Séminaire BOURBAKI\hfill February 1971\\
23rd year, 1970/71, no. 389\\[2.2em]
{\Large\bfseries WORKS OF SHIMURA}\footnote{Text written in March 1971}\\[1.2em]
by Pierre DELIGNE
\end{center}

Let $X/\Gamma$ be the quotient of a symmetric Hermitian domain $X$ by a discrete arithmetic group $\Gamma$ (assumed to be defined by congruence conditions, see 1.7). By Baily and Borel, $X/\Gamma$ is a complex algebraic variety. In the papers [10] to [14], Shimura shows, among other things, for many of these varieties, that they can be defined over number fields which he makes explicit. This is the subject of the present report.

To obtain clean statements, one is forced to use adelic language. This leads one to consider, for a given ``congruence condition'', a scheme over $\C$, a disjoint sum of finitely many varieties of the type $X/\Gamma$ (see §§ 1 and 2). In many cases, this scheme can be defined over a number field $E$, independent of the congruence condition under consideration. Its geometric connected components, however, will be defined over class fields of $E$. This phenomenon, although essential, will be neglected in the rest of this introduction.

In a small number of cases, $X/\Gamma$ can be interpreted as the set of isomorphism classes of complex abelian varieties, endowed with some additional algebraic structures (polarizations, endomorphisms, structures on the points of order $n$). The solution, over $\Q$, of a corresponding moduli problem then provides a scheme $M$ over an explicit number field $F$, whose complex points are $X/\Gamma$. We shall call $M$ a ``model'' of $X/\Gamma$. The fundamental case is that of polarized abelian varieties of the principal series, endowed with a level-$n$ structure (1.6, 1.11, 4.16, 4.17 and 4.21).

In other cases, $X/\Gamma$ can be interpreted as the set of isomorphism classes of complex abelian varieties $A$ endowed with some structures as above, and with some integral cohomology classes of type $(p,p)$
\[
a_i\in H^{2p}(A\times\cdots\times A,\C) \qquad (\text{cf. Mumford [7]}).
\]
In this case, the idea will be to construct a model $M$ of $X/\Gamma$ as a subscheme of a model $M'$ already constructed for a variety $X'/\Gamma'$. Roughly speaking, one first constructs in $M'$ the points of $M$ corresponding to abelian varieties of C.M. type (= maximum of complex multiplication). One then defines $M$ as the closure of the set of these points. This construction rests on the detailed knowledge of abelian varieties of C.M. type provided by [16].

In the cases considered above, one has much information on the fields of definition of the points of $M$ corresponding to abelian varieties of C.M. type. This makes it possible to give partial solutions to Hilbert's twelfth problem (cf. the introduction of [15] and 3.15, 3.16).

Looking at what these cases have in common, one arrives at the notion of a ``canonical model'' (§ 3). One shows that there is at most one ``canonical model'' of the $X/\Gamma$ (5.5). Descent techniques then make it possible to construct some canonical models which seem far from being solutions of moduli problems ([12] [13] [14] and [5]). For the description of these ``strange models'', one is referred to §§ 5 and 6.

We shall denote by $\Z/n(1)$ both the group scheme $\mu_n$ of $n$th roots of unity and the group of $n$th roots of unity in a fixed algebraically closed field $\bar{k}$. For $\bar{k}=\C$, the exponential map identifies $\Z/n(1)$ with $\Z/n$. The $\Z/n(1)$ form a projective system: the transition maps are
\[
\varphi_{n,m}:x\longmapsto x^{n/m}.
\]
We put
\[
\whZ=\limproj_n \Z/n\Z \simeq \prod_p \Z_p,\qquad
\whZ(1)=\limproj_n \Z/n\Z(1),
\]
\[
\Af=\Q\otimes\whZ\simeq \prod_p{}'\Q_p\quad(\text{restricted product}),\qquad
\Af(1)=\Q\otimes\whZ(1)=\Af\otimes_{\whZ}\whZ(1),
\]
and we denote by $\A$ the ring of adeles of $\Q$.

We shall use the following notation:

\noindent $\pii(X)$ (for $X$ a topological space): the set of connected components of $X$. In what follows, $\pii(X)$, endowed with the quotient topology from $X$, will be discrete or compact totally disconnected.

\noindent $G^0$: the connected component of the identity of a topological group $G$ (or of a pointed topological space).

\noindent $G(K)$, $G_K$, $G\otimes_FK$: for $G$ a scheme over $F$ (for instance an algebraic group over $F$) and $K$ an $F$-algebra, $G(K)$ denotes the set of $K$-valued points of $G$, and $G_K$ or $G\otimes_FK$ denotes the scheme over $K$ deduced from $G$ by extension of scalars.

\noindent $E^*$: for $E$ an algebra of finite rank over a field $F$, $E^*$ denotes the algebraic group over $F$ of invertible elements of $E$. For $E$ a number field and $F=\Q$, $E^*(\A)/E^*(\Q)$ is the idele class group of $E$.

We shall say that an algebraic group $G$ over $\Q$ satisfies the Hasse principle if
\[
H^1(\Q,G)=\dfn H^1(\Gal(\overline\Q/\Q),G(\overline\Q))
\]
injects into the product, over all places, of the $H^1(\Q_v,G_{\Q_v})$. We shall use the following theorems.

\paragraph{(0.1) Hasse principle.} A simply connected semisimple group, without a factor of type $E_8$, satisfies the Hasse principle [3].

\paragraph{(0.2)} A simply connected semisimple group $G$ over a non-archimedean local field $K$ satisfies $H^1(K,G)=0$ (a proof independent of the classification is announced in Bruhat--Tits [1]).

\paragraph{(0.3) Strong approximation theorem.} Let $G$ be a simply connected semisimple group over a global field $K$. If $G$ is $K$-simple, and if $v$ is a place of $K$ such that $G(K_v)$ is non-compact, then $G(K_v)G(K)$ is dense in $G(\A)$ (for a proof independent of the classification, see Platonov [9]).

\paragraph{(0.4) Real approximation theorem.} Let $G$ be a connected (linear) algebraic group over $\Q$. Then $G(\Q)$ is dense in $G(\R)$ (to prove this, one reduces easily to the case of tori, which is a corollary of ([4] 5.1)).

\section*{§ 1. The adelic language.}

\paragraph{1.1.} Let $G$ be a reductive group over $\Q$. We denote by $G'$ the derived group of the neutral component of $G$, by $C$ the center of this neutral component, and we put $T=G/G'$. There are exact sequences of algebraic groups
\begin{align}
0&\longrightarrow C\longrightarrow G\longrightarrow G/C\longrightarrow0,\tag{1.1.1}\\
0&\longrightarrow G'\longrightarrow G\xrightarrow{\nu}T\longrightarrow0,\tag{1.1.2}
\end{align}
and the canonical maps from $C$ to $T$ and from $G'$ to $G/C$ are isogenies with kernel $C\cap G'$.

\paragraph{1.2.} Let $\bS$ be the real algebraic group of invertible elements of the $\R$-algebra $\C$. It is also the real algebraic group obtained, by Weil restriction of scalars from $\C$ to $\R$, from the multiplicative group $\Gm$. One has $\bS(\R)=\C^*$, and $\bS_\C\simeq\mathbb G_{m,\C}\times\mathbb G_{m,\C}$. The group $\Hom(\bS_\C,\mathbb G_{m,\C})$ of characters of $\bS_\C$ has as a basis the characters $z$ and $\bar z$ such that the composite
\[
\C^*=\bS(\R)\longrightarrow \bS(\C)\xrightarrow{z,\bar z}\C^*
\]
is respectively the identity and complex conjugation.

Let $V$ be a real vector space. For any representation $h:\bS\to\GL(V)$, one denotes by $V^{pq}$ the subspace of $V_\C$ on which $\bS$ acts by the character $z^p\bar z^q$. The $V^{pq}$ form a bigrading of $V_\C$, and the construction $h\mapsto V^{pq}$ identifies the representations of $\bS$ in $\GL(V)$ with Hodge bigradings of $V_\C$, i.e. with bigradings of $V_\C$ such that $V^{pq}=\overline{V^{qp}}$. We shall denote by $F(h)$ the Hodge filtration of $V_\C$:
\[
F(h)^p=\bigoplus_{p'\ge p}V^{p'q'}.
\]

\paragraph{1.3.} We denote by $r:\mathbb G_{m,\C}\to\bS_\C$ the $\C$-homomorphism such that
\[
(z^p\bar z^q)\circ r=(x\longmapsto x^p),
\]
and by $w:\mathbb G_{m,\R}\to\bS$ the $\R$-homomorphism such that
\[
(z^p\bar z^q)\circ w=(x\longmapsto x^{p+q}).
\]
For a representation $h:\bS\to\GL(V)$ and $v^{pq}\in V^{pq}$, one has
\[
hr(x)\,v^{pq}=x^p v^{pq}\quad\text{and}\quad hw(x)\,v^{pq}=x^{p+q}v^{pq}.
\]
The composite $w:\R^*=\mathbb G_m(\R)\xrightarrow{w}\bS(\R)=\C^*$ is the evident inclusion.

\paragraph{1.4.} Let $h:\bS\to G_\R$ be a homomorphism. For any representation $\rho:G\to\GL(V)$ of $G$ on a $\Q$-vector space $V$, $\rho h$ defines a bigrading of $V_\C$ (1.2). If $V$ is faithful, and if $G$ is the subgroup of $\GL(V)$ fixing some tensors $s_i$, the construction $h\mapsto(V^{pq})$ identifies the homomorphisms $h$ with Hodge bigradings of $V_\C$ such that the $s_i$ are of type $(0,0)$.

\paragraph{1.5.} Let $h_o:\bS\to G_\R$ be a homomorphism satisfying the following conditions (cf. [7] and [2]; (1.5.2) is motivated by Griffiths' transversality theorem).

\begin{enumerate}[label={(1.5.\arabic*)},leftmargin=3.0em]
\item The image of $h_ow:\Gm\to G_\R$ is central. We shall call $h_ow$ the weight of $h_o$;
\item the Hodge bigrading of the complexification $\Lie(G)_\C$ of the space of the adjoint representation (1.2) is of type $(-1,1)+(0,0)+(1,-1)$;
\item the automorphism $\ad h_o(i)$ of $G$, an involution by (1.5.1), induces a Cartan involution of $G'$.
\end{enumerate}

Let $K_\infty$ be the centralizer of $h_o$ in $G(\R)$; it contains the center of $G(\R)$, and $K_\infty\cap G'(\R)^0$ is a maximal compact subgroup of $G'(\R)^0$, identical with the centralizer in $G'(\R)^0$ of $h_o(i)$. The space $K_\infty\backslash G(\R)$ is identified with the set $X$ of conjugates of $h_o$ by an element of $G(\R)$. One checks that it has one and only one complex structure such that, for every representation $\rho:G\to\GL(V)$, the filtration $F(\rho h)$ of $V_\C$ depends holomorphically on $h\in X$. This structure is invariant on the right under $G(\R)$, and the connected components of $X$ are symmetric Hermitian domains. We shall denote by $X^0$ the connected component of $h_o$.

\paragraph{1.6. Example I.} Let $\Gp(V)$ be the group of symplectic similitudes of a real vector space $V$ endowed with a non-degenerate alternating form $\psi$. There is then one and only one conjugacy class $X$ of homomorphisms $h:\bS\to\Gp(V)$ satisfying conditions (1.5.1) and of weight $-1$, i.e. of weight
\[
hw:\Gm\to\Gp(V):x\longmapsto\text{homothety of ratio }x^{-1}.
\]
The space $X$ is identified with the sum of two Siegel half-spaces.

Construction 1.5 identifies the $h\in X$ with decompositions
\[
V_\C=V^{-1,0}\oplus V^{0,-1}
\]
of $V_\C=V\otimes_\R\C$, such that $V^{pq}=\overline{V^{qp}}$, such that $V^{0,-1}$ is totally isotropic for $\psi$, and such that, if $C$ denotes the endomorphism of $V$ which induces multiplication by $i^{p-q}$ on $V^{pq}$, the symmetric form $\psi(x,Cy)=\psi(x,h(i)y)$ on $V$ is definite ($>0$ or $<0$).

\paragraph{1.7.} Let $\Gamma$ be an arithmetic subgroup of $G(\Q)$. We are interested in the case where $\Gamma$ is defined by congruence conditions. If $V_\Z$ is an integral lattice in a faithful representation $V$ of $G$, this means that there exists $n$ such that $\Gamma$ contains, with finite index, the group $\Gamma_n\subset G(\Q)$ of $\gamma\in G(\Q)$ such that $\gamma\in G(\R)^0$, $\gamma V_\Z=V_\Z$, and $\gamma\equiv\id\pmod n$. By Baily and Borel, the quotient $X^0/\Gamma$ is endowed with a natural structure of quasi-projective scheme over $\C$. Let $K_n$ be the compact open subgroup of $G(\Af)$ formed by the $k=(k_p)$ such that
\[
k_p\cdot V_\Z\otimes\Z_p=V_\Z\otimes\Z_p
\quad\text{and}\quad
k_p\equiv1\pmod n\quad\text{for }p\mid n.
\]
The group $\Gamma_n$ is then the intersection, inside $G(\A)$, of $G(\Q)$ and $G(\R)^0\times K_n$.

\paragraph{1.8.} Let $K$ be a compact open subgroup of $G(\Af)$. The space
\[
K_\infty\times K\backslash G(\A)/G(\Q)
=K\backslash X\times G(\Af)/G(\Q)
=X\times(K\backslash G(\Af))/G(\Q)
\]
is then a sum of a finite family of spaces of the type considered in 1.7. It is therefore a quasi-projective scheme over $\C$. We shall denote it by ${}_K\MC(G,h_o)$ (or simply ${}_K\MC(G)$ or ${}_K\MC$).

For $K$ becoming smaller and smaller, these schemes form a projective system of schemes with finite and surjective transition morphisms
\[
{}_K\MC(G,h_o)\qquad(K\to0). \tag{1.8.1}
\]
We shall denote by $\MC(G,h_o)$ (or simply $\MC(G)$, or $\MC$) the quasi-compact separated scheme (not of finite type for $G\ne\{1\}$) which is the projective limit of (1.8.1)
\[
\MC(G,h_o)=\limproj_K K_\infty\times K\backslash G(\A)/G(\Q). \tag{1.8.2}
\]
It must be considered as an avatar of the projective system (1.8.1). The group $G(\Af)$ acts on $\MC$, or, if one prefers, on the pro-object defined by the projective system (1.8.1), but of course not on the individual ${}_K\MC$ (cf. 3.2). One has
\[
K\backslash\MC={}_K\MC,
\]
so that
\[
\MC=\limproj_K K\backslash\MC\qquad(K\text{ compact open in }G(\Af)). \tag{1.8.3}
\]
This is expressed by saying that $G(\Af)$ acts continuously on $\MC$.

The following construction will be useful for interpreting the ${}_K\MC$ as moduli schemes.

\paragraph{1.9.} Let $H$ be an algebraic group over $\Q$ and $V$ a faithful representation of $H$. We define an $H$-structure $\bar\iota$ on a vector space $W$ to be an isomorphism from $V$ to $W$, given modulo $H(\Q)$: $\bar\iota\in\Isom(V,W)/H(\Q)$. If $H(\Q)$ is the group of automorphisms of a structure $s$ of species $\Sigma$ on $V$, then an $H$-structure on $W$ is identified with a structure $t$ of the same species on $W$, isomorphic to $s$ (with $\bar\iota=\Isom((V,s),(W,t))$). Let $W$ be endowed with an $H$-structure $\bar\iota$. The algebraic group $\iota H\iota^{-1}\subset\GL(W)$ ($\iota\in\bar\iota$) depends only on $\bar\iota$. It is denoted $H^{\bar\iota}$. Let $A$ be a $\Q$-algebra. The set $\Isom_{\bar\iota}(W\otimes A,V\otimes A)$ of permissible isomorphisms from $W\otimes A$ to $V\otimes A$ is the set of isomorphisms $k$ such that, for $\iota\in\bar\iota$, one has $k\iota\in H(A)$.

\paragraph{1.10.} Let $(G,h_o)$ be as in 1.5, let $V$ be a faithful linear representation of $G$, let $K$ be a compact open subgroup of $G(\Af)$, and consider objects $H$ of the following type: $H$ consists of
\begin{enumerate}[label=(\alph*),leftmargin=2.6em]
\item a vector space $H_\Q$ over $\Q$, endowed with a $G$-structure $\bar\iota$ (1.9);
\item a homomorphism $h:\bS\to G^{\bar\iota}_\R$ such that, for $\iota\in\bar\iota$, $\iota^{-1}h\iota:\bS\to G_\R$ is conjugate to $h_o$;
\item a class $\bar k\in K\backslash\Isom_{\bar\iota}(H\otimes\Af,V\otimes\Af)$ (1.9).
\end{enumerate}

\paragraph{1.11. Example I (continued from 1.6).} Let $V$ be a $\Q$-vector space endowed with a non-degenerate alternating form $\psi$, let $G$ be the group of symplectic similitudes of $V$, let $V_\Z$ be an integral lattice on which $\psi$ has integral values and discriminant $1$, and let $h_o$ be as in 1.6. Take
\[
K=\prod_\ell \Gp(V_\Z\otimes\Z_\ell).
\]
A datum (a) is interpreted as a vector space $H$ of the same dimension as $V$, endowed with an alternating form $\psi$ given up to a rational factor. A datum (b) is interpreted as a Hodge bigrading of $H_\C$, of type 1.6. A datum (c) is interpreted as an integral lattice $H_\Z$ of $H$, on which a rational multiple of $\psi$ has integral values and discriminant $1$: $\bar k$ is the set of $k$ such that, for every $\ell$, $V_\Z\otimes\Z_\ell=k_\ell(H_\Z\otimes\Z_\ell)$. One can normalize $\psi$ by the conditions of having integral values and discriminant $1$ on $H_\Z$, and of satisfying
\[
\psi(x,Cx)>0 \qquad (\text{cf. }1.6).
\]
Let $n$ be an integer, and let $K_n$ be the subgroup of $K$ formed by the $k$ such that $k\equiv1\pmod n$. For $K_n$, a datum (c) is interpreted as $H_\Z$ as above, plus a symplectic similitude
\[
H_\Z/nH_\Z\xrightarrow{\sim}V_\Z/nV_\Z.
\]

\paragraph{1.12.} Under the hypotheses of 1.10, the points of ${}_K\MC(G,h_o)$ are in bijective correspondence with the isomorphism classes of objects $H$ of type 1.10. Indeed, let
\[
H=(H_\Q,\bar\iota,h,\bar k).
\]
To $H$ endowed with $\iota\in\bar\iota$, one associates the morphism $\iota^{-1}h\iota:\bS\to G_\R$, an element of $X\simeq K_\infty\backslash G(\R)$, and $\bar k$ in $K\backslash G(\Af)$, hence in total an element of $K\backslash X\times G(\Af)$; to $H$ alone, one associates only an element of
\[
{}_K\MC(G,h_o)=K\backslash X\times G(\Af)/G(\Q)=K_\infty\times K\backslash G(\A)/G(\Q),
\]
which uniquely determines the isomorphism class of $H$.

In certain cases (see 4.11), to give an object $H$ as above amounts to giving an abelian variety endowed with some supplementary structures. One then obtains an interpretation of ${}_K\MC(G)$ (and of $\MC(G)$) as the coarse moduli scheme for this type of abelian variety.

\paragraph{1.13.} Let $G^i$ ($i=1,2$) be two reductive groups, let $h^i:\bS\to G^i_\R$ be homomorphisms satisfying the conditions of 1.5, and let $X^i$ be the set of conjugates of $h^i$ by an element of $G^i(\R)$. Let $G=G^1\times G^2$, let $h:\bS\to G_\R$ be the morphism $(h^1,h^2)$, and let $X=X^1\times X^2$ be the set of conjugates of $h$. In the evident way one defines an isomorphism
\[
\MC(G^1,h^1)\times\MC(G^2,h^2)=\MC(G,h). \tag{1.13.1}
\]
For $K^i$ compact open in $G^i(\Af)$ and $K=K^1\times K^2$, one likewise has
\[
{}_{K^1}\MC(G^1,h^1)\times{}_{K^2}\MC(G^2,h^2)={}_K\MC(G,h). \tag{1.13.2}
\]

\paragraph{1.14.} Let $G^1,G^2,h^1,h^2$ be as in 1.13. We shall denote by $u:(G^1,h^1)\to(G^2,h^2)$ a homomorphism $u:G^1\to G^2$ such that $uX^1\subset X^2$. Such a homomorphism defines
\[
u:\MC(G^1,h^1)\longrightarrow\MC(G^2,h^2). \tag{1.14.1}
\]
For $K^1$ compact open in $G^1(\Af)$, with $u(K^1)\subset K^2$, it defines
\[
u(K^1,K^2):{}_{K^1}\MC(G^1,h^1)\longrightarrow{}_{K^2}\MC(G^2,h^2). \tag{1.14.2}
\]

\paragraph{Proposition 1.15.} With the notation of 1.14, suppose that $G^1$ is a subgroup of $G^2$. Then, for every compact open subgroup $K^1$ of $G^1(\Af)$, there exists a compact open subgroup $K^2\supset K^1$ of $G^2(\Af)$ such that $u(K^1,K^2)$ identifies ${}_{K^1}\MC(G^1,h^1)$ with a closed subscheme of ${}_{K^2}\MC(G^2,h^2)$.

If the $K^i$ are small enough that the map from $X^i\times(K^i\backslash G^i(\Af))$ to ${}_{K^i}\MC(G^i,h^i)$ is étale, then the map $u(K^1,K^2)$ is finite and unramified. It suffices to treat this case, and to show that $u(K^1,K^2)$ is injective for $K^2\supset K^1$ sufficiently small. The graph of the equivalence relation
\[
u(K^1,K^2)(x)=u(K^1,K^2)(y)
\]
is a closed subscheme of $({}_{K^1}\MC(G^1,h^1))^2$, decreasing with $K^2$, hence stationary for $K^2\supset K^1$ sufficiently small. It is enough to show that
\[
u(K^1):{}_{K^1}\MC(G^1,h^1)\longrightarrow\limproj_{K^2\supset K^1}{}_{K^2}\MC(G^2,h^2)=\dfn{}_{K^1}\MC(G^2,h^2)
\]
is injective. One has
\[
{}_{K^1}\MC(G^1,h^1)=K^1\backslash \MC(G^1,h^1),
\qquad
{}_{K^1}\MC(G^2,h^2)=K^1\backslash \MC(G^2,h^2),
\]
so it is enough to prove the

\paragraph{Variant 1.15.1.} $u:\MC(G^1,h^1)\longrightarrow\MC(G^2,h^2)$ is injective.



This assertion follows easily from the following two lemmas.

\paragraph{Lemma 1.15.2.} Let $U^i\subset C^i(\Q)$ be the group of units of the centre of $G^i$, and put
\[
G^i(\Q)^\wedge=\limproj_n\bigl(G^i(\Q)/(U^i)^n\bigr).
\]
Then
\[
\MC(G^i,h^i)=K^i_\infty\backslash G^i(\A)/G^i(\Q)^\wedge .
\]
This is a corollary of Chevalley's theorem, according to which the $(U^i)^n$ are congruence subgroups of $U^i$.

\paragraph{Lemma 1.15.3.} The map
\[
u:G^1(\Af)/G^1(\Q)^\wedge\longrightarrow G^2(\Af)/G^2(\Q)^\wedge
\]
is injective.

By left translation it is enough to show that $u^{-1}(\{e\})=\{e\}$. Consider the diagram
\[
\begin{array}{ccccccccc}
0&\longrightarrow&G^1(\Q)^\wedge&\longrightarrow&G^2(\Q)^\wedge&\longrightarrow&(G^2/G^1)(\Q)^\wedge\\
&&\downarrow&&\downarrow&&\downarrow\!\\
0&\longrightarrow&G^1(\Af)&\longrightarrow&G^2(\Af)&\longrightarrow&(G^2/G^1)(\Af)\\
&&\downarrow&&\downarrow&&\\
&&G^1(\Af)/G^1(\Q)^\wedge&\longrightarrow&G^2(\Af)/G^2(\Q)^\wedge.&&
\end{array}
\]
where
\[
(G^2/G^1)(\Q)^\wedge=\limproj_n (G^2/G^1)(\Q)/(U^2/U^1)^n .
\]

\section*{§ 2. Connected components.}

\paragraph{2.1.} Let $G$ be a connected reductive group over $\Q$. We use the notation of 1.1 and assume that $G'$ admits no non-trivial invariant subgroup $H$ (defined over $\Q$) such that $H(\R)$ is compact. This hypothesis does not exclude the possibility that $G'_\R$ has compact factors. We denote by $\widetilde G'$ the simply connected covering of the derived group $G'$.

\paragraph{Proposition 2.2.} (i) The group $G(\Af)$ acts transitively on the set
\[
\pi_0(G(\A)/G(\Q))
\]
of connected components of $G(\A)/G(\Q)$.

(ii) The group $\widetilde G'(\A)$ acts trivially on $\pi_0(G(\A)/G(\Q))$.

(iii) The quotient of $G(\A)$ by the image of $\widetilde G'(\A)$ is a commutative group.

Assertion (i) follows from (0.4):
\[
G(\A)=G(\Af)G(\R)^0G(\Q).
\]
Since $\widetilde G'$ is simply connected, $\widetilde G'(\R)$ is connected. By (0.3), $\widetilde G'(\R)\widetilde G'(\Q)$ is dense in $\widetilde G'(\A)$, so that $\widetilde G'(\A)/\widetilde G'(\Q)$ is connected. For $x\in G(\A)$ one has $\widetilde G'(\A)x=x\widetilde G'(\A)$; hence the image of $\widetilde G'(\A)x$ in $G(\A)/G(\Q)$ is an image of $\widetilde G'(\A)/\widetilde G'(\Q)$: it is connected, and (ii) follows.

Let $Z$ be the centre of $\widetilde G'$. Assertion (iii) follows from the fact that the commutator map
\[
(x,y)\longmapsto xyx^{-1}y^{-1}:G\times G\longrightarrow G
\]
has a factorization
\[
G\times G\longrightarrow G/C\times G/C\longleftarrow \widetilde G'/Z\times \widetilde G'/Z\longrightarrow \widetilde G'\longrightarrow G.
\]

\paragraph{Definition 2.3.} We denote by $\pi(G)$ the commutative quotient of $G(\A)$, or of $G(\Af)$, which acts faithfully on $\pi_0(G(\A)/G(\Q))$.

In other words, $\pi(G)$ is $\pi_0(G(\A)/G(\Q))$, endowed with its ``evident'' commutative group structure. We recall the following theorem.

\paragraph{Theorem 2.4.} Under the hypotheses of 2.1, if $G'$ is simply connected, then the canonical homomorphism
\[
\pi(G)=\pi_0(G(\A)/G(\Q))\longrightarrow \pi_0(T(\A)/T(\Q)) \tag{2.4.1}
\]
is bijective.

To prove the surjectivity of (2.4.1), by 2.2(i) applied to $T$ it is enough to show that $G(\Af)$ maps onto $T(\Af)$. Since the kernel $G'$ of $\nu:G\to T$ is connected, one knows that, for almost all $p$, $\nu(G(\Z_p))=T(\Z_p)$; this formula has a meaning for almost all $p$, and is a corollary of Lang's theorem applied to $G'$ reduced modulo $p$. It is therefore enough to prove that, for every $p$,
\[
\nu(G(\Q_p))=T(\Q_p),
\]
which follows from (0.2) applied to $G'$.

Theorem 2.4 is equivalent to the following more concrete statement.

\paragraph{Variant 2.5.} For every compact open subgroup $K$ of $G(\Af)$, the map
\[
\nu:\pi_0(K\backslash G(\A)/G(\Q))\longrightarrow\pi_0(\nu(K)\backslash T(\A)/T(\Q)) \tag{2.5.1}
\]
is bijective.

Indeed,
\[
\pi_0(G(\A)/G(\Q))=\limproj_K\pi_0(K\backslash G(\A)/G(\Q)),
\]
and similarly for $T$. Also,
\[
\pi_0(K\backslash G(\A)/G(\Q))=G(\R)^0\backslash K\backslash G(\A)/G(\Q).
\]
We show that if $x,y\in G(\A)$ have the same image in the right-hand side of (2.5.1), then they have the same image in the left-hand side. Left translation by $y^{-1}$, which replaces $K$ by $y^{-1}Ky$, allows us to suppose that $y=e$. Modifying $x$ by an element of $G(\R)^0\times K$, we may suppose that $\nu(x)=\tau$ with $\tau\in T(\Q)$. By the Hasse principle (0.1) for $G'$, the equation $\nu(y)=\tau$ has a solution in $G(\Q)$. Correcting $x$ by $y$, we may suppose that $x\in G'(\A)$. By 2.2(ii), the image of $x$ in $G(\A)/G(\Q)$ is in the neutral component, and a fortiori $x$ and $y$ have the same image in the left-hand side of 2.5.1.

This proof seems to use, via (0.1), that $G'$ has no factor of type $E_8$. In fact, the class in $H^1(\Q,G')$ which obstructs the equation $\nu(y)=\tau$ comes from the centre of $G'$, and the $E_8$ factors, being adjoint, count for nothing.

\paragraph{Proposition 2.6.} One has
\[
K_\infty=C(\R)\bigl(K_\infty\cap G'(\R)^0\bigr).
\]
The group $K_\infty\cap G'(\R)^0$ is connected, and
\[
\pi_0(K_\infty)\xrightarrow{\sim}\operatorname{Im}\bigl(\pi_0(C(\R))\longrightarrow\pi_0(G(\R))\bigr).
\]

Let $h':\bS\to G_\R\to(G/C)_\R$. The centralizer of $h'$ in $(G/C)(\R)$ is connected, because it is a compact group whose complexification is connected, as the centralizer of a torus. In particular, the image of $K_\infty$ in $G/C(\R)$ lies in the neutral component, the image of $G'(\R)^0$, and
\[
K_\infty=C(\R)\bigl(K_\infty\cap G'(\R)^0\bigr).
\]
The group $K_\infty\cap G'(\R)^0$ is connected as a maximal compact subgroup of a connected group. This proves 2.6.

\paragraph{2.7.} The image of $\pi_0(K_\infty)$ in $\pi_0(G(\A)/G(\Q))$ depends only on the conjugacy class of $h$, and
\[
\pi_0(\MC(G,h))=\limproj_K\pi_0(K_\infty\times K\backslash G(\A)/G(\Q))
\]
is identified with $\pi(G)/\pi_0(K_\infty)$.

In particular, if $G'$ is simply connected, one has
\[
\pi_0(\MC(G,h))\xrightarrow{\sim}\pi_0(T(\A)/T(\Q))/\pi_0(K_\infty). \tag{2.7.1}
\]
More concretely, for every compact open subgroup $K$ of $G(\Af)$, $\pi_0({}_K\MC(G,h))$ is then a principal homogeneous space under
\[
\nu(K_\infty\times K)\backslash T(\A)/T(\Q).
\]

\section*{§ 3. Models.}

Let $G$ be a reductive group over $\Q$, let $h:\bS\to G_\R$ satisfy the conditions of 1.5, and let $E$ be a field endowed with a complex embedding $\rho:E\to\C$. Put $\MC(G)=\MC(G,h)$.

\paragraph{Definition 3.1.} A model over $E$ of $\MC(G)$ consists of
\begin{enumerate}[label=(\alph*),leftmargin=2.4em]
\item a scheme $M$ over $E$, endowed with a continuous action of $G(\Af)$;
\item an isomorphism $m$ from $M\otimes_{E,\rho}\C$ to $\MC(G)$, compatible with the action of $G(\Af)$.
\end{enumerate}
Let $F$ be a finite extension of $E$, endowed with a complex embedding extending that of $E$. If $M_E(G,h)$ is a model of $\MC(G,h)$ over $E$, we denote by $M_F(G,h)$ the model $M_E(G,h)\otimes_E F$ of $\MC(G,h)$ over $F$.

\paragraph{Remark 3.2.} Giving a scheme $M$ over $E$, endowed with a continuous action of $G(\Af)$, is equivalent to giving
\begin{enumerate}[label=(\alph*),leftmargin=2.4em]
\item for every compact open subgroup $K$ of $G(\Af)$, a scheme ${}_K M$ over $E$;
\item for compact open subgroups $K,L$ of $G(\Af)$ and for $x\in G(\Af)$ such that $xKx^{-1}\subset L$, a homomorphism
\[
J_{L,K}(x):{}_K M\longrightarrow {}_L M,
\]
these homomorphisms satisfying
\begin{align*}
(\alpha)\quad&J_{M,L}(y)J_{L,K}(x)=J_{M,K}(yx),\tag{3.2$\alpha$}\\\
(\beta)\quad&J_{K,K}(x)=\id\qquad\text{if }x\in K,\tag{3.2$\beta$}\\\
(\gamma)\quad&\text{for }K\text{ normal in }L,\; J\text{ defines an action of }L/K\text{ on }{}_K M,\\
&\text{and }J_{L,K}(e)\text{ defines }(L/K)\backslash {}_K M\xrightarrow{\sim}{}_L M.\tag{3.2$\gamma$}
\end{align*}
\end{enumerate}
One has
\[
M=\limproj_K {}_K M,\tag{3.2.1}
\]
and
\[
{}_K M=K\backslash M.\tag{3.2.2}
\]

\paragraph{3.3.} Let $M(G)$ be a model over $E$ of $\MC(G)$. For $K$ compact open in $G(\Af)$, put ${}_K M(G)=K\backslash M(G)$. The ${}_K M(G)$ are quasi-projective schemes over $E$, not necessarily geometrically connected. The structural isomorphism $m$ induces isomorphisms
\[
{}_K M(G)\otimes_{E,\rho}\C\xrightarrow{\sim}{}_K\MC(G).\tag{3.3.1}
\]

\paragraph{3.4.} Suppose that $G$ satisfies the hypotheses of 2.1, and let $(M,m)$ be a model of $\MC(G,h)$ over $E$. Let $\overline E$ be the algebraic closure of $E$ in $\C$. Since $M$ is defined over $E$, the Galois group $\Gal(\overline E/E)$ acts on the profinite set
\[
\pi_0(M\otimes_E\overline E)=\limproj_K\pi_0({}_K M\otimes_E\overline E)\xrightarrow{\sim}_{m}\pi_0(\MC(G,h)).
\]
The group $G(\Af)$ acts on $\pi_0(M\otimes_E\overline E)$ through its quotient $\pi(G)/\pi_0(K_\infty)$, and $\pi_0(M\otimes_E\overline E)$ is a principal homogeneous space under the commutative group $\pi(G)/\pi_0(K_\infty)$ (2.6). The action of $\Gal(\overline E/E)$ commutes with that of $G(\Af)$, hence with that of $\pi(G)/\pi_0(K_\infty)$. The action of an element $\sigma$ of $\Gal(\overline E/E)$ can therefore only be the translation defined by an element $\lambda(\sigma)$ of $\pi(G)/\pi_0(K_\infty)$, and $\lambda$ is a homomorphism
\[
\lambda_M:\Gal(\overline E/E)\longrightarrow \pi(G)/\pi_0(K_\infty).\tag{3.4.1}
\]
Assume that $E$ is a number field. Class field theory identifies the largest abelian quotient of $\Gal(\overline E/E)$ with $\pi_0(E^*(\A)/E^*(\Q))$, and identifies (3.4.1) with
\[
\lambda_M:\Gal(\overline E/E)^{\ab}=\pi_0(E^*(\A)/E^*(\Q))\longrightarrow\pi(G)/\pi_0(K_\infty).\tag{3.4.2}
\]
In the special case where $G'$ is simply connected, the homomorphism (3.4.2) is identified, by (2.7.1), with
\[
\lambda_M:\pi_0(E^*(\A)/E^*(\Q))\longrightarrow\pi_0(T(\A)/T(\Q))/\pi_0(K_\infty).\tag{3.4.3}
\]

\paragraph{Definition 3.5.} The homomorphisms (3.4.1), (3.4.2), and (3.4.3) are called the reciprocity law of the model $M$. If $G'$ is simply connected, and if $\lambda_M$ comes from a homomorphism of algebraic groups over $\Q$ from $E^*$ to $T$, the latter will also be called the reciprocity law of $M$.

\paragraph{3.6.} There is a rule, applicable to each of the models constructed by Shimura, for determining $E$ and the reciprocity law $\lambda_M$ from $G$ and $h$. The field $E$, for arbitrary $G$, and $\lambda_M$, for simply connected $G'$, are described as follows.

\paragraph{3.7.} The composite homomorphism $hr$ (1.3)
\[
hr:\mathbb G_{m,\C}\xrightarrow{r}\bS_\C\xrightarrow{h}G_\C
\]
is a homomorphism over $\C$ between algebraic groups defined over $\Q$. The subfield $E$ of $\C$ is the field of definition $E(G,h)$ of the conjugacy class of this homomorphism.

For $G$ semisimple adjoint, the field $E(G,h)$ can be described as follows. Let $\Delta$ be the Dynkin diagram of $G_\C$, and let $|\Delta|$ be the set of vertices of $\Delta$. The Galois group $\Gal(\C/\Q)$ acts on $\Delta$.

Let $H$ be a maximal torus of $G_\C$, endowed with a system of simple roots $(\alpha_i)_{i\in |\Delta|}$. The $\alpha_i$ identify $H$ with $\Gm^{\Delta}$. For every homomorphism $u:\mathbb G_{m,\C}\to G_\C$, there is then a unique homomorphism $u':\mathbb G_{m,\C}\to H=\mathbb G_{m,\C}^{\Delta}$,
\[
x\longmapsto (x^{n_i})_{i\in |\Delta|}\qquad(n_i\ge0),
\]
which is conjugate to $u$. The construction $u\mapsto \eta(u)=(n_i)_{i\in |\Delta|}$ identifies the set of conjugacy classes of maps from $\mathbb G_{m,\C}$ to $G_\C$ with the set of functions from $|\Delta|$ to $\mathbb N$. It follows that $E(G,h)$ is defined by the subgroup of $\Gal(\C/\Q)$ which stabilizes $\eta(hr)$.

\paragraph{Proposition 3.8.} Let
\[
G/C=\prod_{i=1}^g G_i
\]
be the decomposition of the adjoint group $G/C$ into $\Q$-simple factors, and let $h_i$ be the composite
\[
h_i:\bS\xrightarrow{h}G_\R\longrightarrow (G_i)_\R.
\]
(i) The subfield $E(G,h)$ of $\C$ is the compositum of the subfields $E(T,\nu h)$ and $E(G_i,h_i)$ ($1\le i\le g$) of $\C$.

(ii) If the opposition involution of the Dynkin diagram of $G_i$ respects $\eta(h_ir)$, then $E(G_i,h_i)$ is totally real. Otherwise, $E(G_i,h_i)$ is a totally imaginary quadratic extension of a totally real number field.

Assertion (i) is easy to check, and it is enough to prove (ii) for $G$ $\Q$-simple and adjoint. Let $\Delta$ be its Dynkin diagram. The hypotheses 1.5 imply that $G_\R$ has a compact maximal torus; this implies that complex conjugation acts on $\Delta$ by the opposition involution $\iota$. This involution is in the centre of the automorphism group of $\Delta$.

Let $E'$ be the Galois extension of $\Q$ defined by the subgroup of $\Gal(\C/\Q)$ which acts trivially on $\Delta$. Then the image $c$ of complex conjugation is central in $\Gal(E'/\Q)$; hence $E'$ is either totally real or a totally imaginary quadratic extension of a totally real field. Finally, $c$ is the identity on $E(G,h)\subset E'$ if and only if $\iota$ respects $\eta(hr)$, and the assertion follows.

\paragraph{3.9.} We use the notation of 3.6. Suppose that $G'$ is simply connected, that $E$ contains $E(G,h)$, and that the complex embedding of $E$ induces that of $E(G,h)$. The composite morphism
\[
r''(h):\mathbb G_{m,\C}\xrightarrow{r}\bS_\C\xrightarrow{h}G_\C\xrightarrow{\nu}T_\C
\]
depends only on the conjugacy class of $hr$. It is therefore defined over $E$, i.e. it is obtained by extension of scalars from $E$ to $\C$ from a homomorphism of algebraic groups over $E$, again denoted
\[
r''(h):\Gm\longrightarrow T_E.
\]
Apply Weil restriction of scalars, and let $r'(h):E^*\to T$ be the composite
\[
E^*=\Res_{E/\Q}\Gm\xrightarrow{\Res_{E/\Q}(r''(h))}\Res_{E/\Q}(T_E)\xrightarrow{N_{E/\Q}}T.
\]
The reciprocity law of $M$ is the inverse $r(G,h)$ of $r'(h)$:
\[
r(G,h)=r'(h)^{-1}:E^*\longrightarrow T .\tag{3.9.1}
\]

\paragraph{3.10. Example II.} Let $H$ be a torus and let $h:\bS\to H_\R$. The conditions of 1.5 are automatically satisfied. The ${}_K\MC(H,h)$ are finite sets. A finite reduced scheme over the field $E(H,h)$ (3.7) is the same as a Galois set. It follows that, up to unique isomorphism, there is a unique model of $\MC(H,h)$ over $E(H,h)$, with reciprocity law given by 3.9.1. It is denoted $M(H,h)$.

\paragraph{3.11.} Let $F$ be a finite extension of $E$, endowed with a complex embedding extending that of $E$. The diagram
\[
\begin{array}{ccc}
F^*&\xrightarrow{N_{F/E}}&E^*\\
&\searrow&\downarrow\\
&&T
\end{array}
\]
is commutative. If $M_E(G,h)$ is a model over $E$ of $\MC(G,h)$ whose reciprocity law is (3.9.1), it follows that $M_F(G,h)$ (3.1) again has (3.9.1) as its reciprocity law.

\paragraph{3.12.} Let $u:(G_1,h_1)\to(G_2,h_2)$ be as in 1.14, and let $M^i(G_i,h_i)$ be a model of $\MC(G_i,h_i)$ over $E_i$. Let $E$ be the compositum of the subfields $E_1$ and $E_2$ of $\C$. With the notation of 3.1 one has
\[
M^i(G_i,h_i)\otimes_{E_i}E=M_E(G_i,h_i)\qquad(i=1,2).
\]
It therefore makes sense to ask whether
\[
u:\MC(G_1,h_1)\longrightarrow\MC(G_2,h_2)
\]
is defined over $E$.

\paragraph{Definition 3.13.} Let $G$ be a connected reductive group over $\Q$ satisfying the condition of 2.1, let $h:\bS\to G_\R$ be a homomorphism satisfying the conditions of 1.5, and let $E$ be an extension of $E(G,h)$ endowed with a complex embedding extending that of $E(G,h)$. A model $M_E(G,h)$ of $\MC(G,h)$ over $E$ is called weakly canonical if, for every torus $u:H\to G$ in $G$, endowed with $h':\bS\to H_\R$ such that $uh'$ is the conjugate of $h$ by an element of $G(\R)$, the morphism (1.14)
\[
\MC(H,h')\longrightarrow\MC(G,h)
\]
is defined, in the sense of (3.10) and (3.12), over the compositum $E(H,h')\cdot E\subset\C$. A model $M_E(G,h)$ is called canonical if it is weakly canonical and if $E=E(G,h)$.

\paragraph{3.14. Translation.} Let $M$ be a canonical model of $\MC(G,h)$, with reciprocity law $\lambda_M$. If $G'$ is simply connected, then $\lambda_M$ is given by 3.9.1 (5.6). For every compact open subgroup $K$ of $G(\Af)$, again denote by $\nu(K)$ the image of $K$ in $\pi(G)/\pi_0(K_\infty)$, and let $E(K)$ be the extension of $E(G,h)$ defined by the idèle class group $\lambda_M^{-1}(\nu(K))$. One verifies that $E(K)$ is the field of definition of the neutral component ${}_K\MC^0$ of ${}_K\MC$, considered as a subscheme of ${}_K\MC={}_K M\otimes_E\C$.

By definition, ${}_K\MC^0$ is obtained by extension of scalars from $E(K)$ to $\C$ from a subscheme ${}_K M'$ of ${}_K M\otimes_EE(K)$. The scheme ${}_K M'/E(K)$ is geometrically connected over $E(K)$, and one can also describe ${}_K M'$ as the union of those connected components, over $E$, of ${}_K M$ which, after extension of scalars from $E$ to $\C$, contain the origin:
\[
\begin{array}{ccc}
{}_K\MC^0&\subset&{}_K\MC\\
\downarrow&&\downarrow\\
\Spec E(K)&\longrightarrow&\Spec\C\\
\downarrow&&\\
\Spec E&=&\Spec E.
\end{array}
\]
Shimura usually states his results in terms of the system of varieties ${}_K M'/E(K)$ and the homomorphisms of type 2.2 between them. For a typical statement, see [14, p. 146]. For the translation, cf. [14, 2.7].




\paragraph{3.15.} For $u:(H,h')\longrightarrow(G,h)$ as in 3.13, the image of $M_{\C}(H,h')$ in ${}_K M_{\C}(G,h)$ is a finite set. The hypothesis is that this part of ${}_K M_{\C}(G,h)$ be defined over $E(H,h')$, that its points be defined over an abelian extension of $E(H,h')$, and that $\Gal(\overline{E(H,h')}/E(H,h'))$ permute them in a prescribed manner.

The points of ${}_K M_{\C}(G,h)$ lying in the image of some $M_{\C}(H,h')$ will be called \emph{special}.

Suppose that a projective embedding
\[
q:{}_K M_{\C}^{0}\longrightarrow \mathbb P^r
\]
is known, and is defined over $E(K)$. It may be interpreted as
\[
\widetilde q:X^0\longrightarrow \mathbb P^r
\]
(cf. 1.7, 1.8). If $uh'\in X$ lies in $X^0$, then the coordinates of $\widetilde q(uh')$ generate, over $E(K)E(H,h')$, an abelian extension of $E(H,h')$ which can be made explicit as a class field.

\paragraph{3.16.} The classical example of such a situation is provided by the moduli scheme of elliptic curves, corresponding to $G=\GL_2$, $K=\GL_2(\widehat{\Z})$, $h$ as in 1.6, and by those of its points defined by elliptic curves with complex multiplication.

Let $j$ be the modular invariant, viewed as a function on the Poincaré half-plane $X^0$. If $\tau\in X^0$ belongs to an imaginary quadratic field $K$, put
\[
L(\tau)=\Z\oplus \Z\tau\subset\C,
\qquad
\mathcal O(\tau)=\{\,k\in K\mid kL(\tau)\subset L(\tau)\,\}.
\]
The field $K(j(\tau))$ is the abelian extension of $K$ whose Galois group is the ideal class group of $\mathcal O(\tau)$. If $\sigma$ in Galois corresponds to an invertible module $L$, and if
\[
L(\tau')\simeq L(\tau)\otimes_{\mathcal O(\tau)}L,
\]
then $j(\tau)=j(\tau')^\sigma$.

For more unusual explicit examples, see [15], pp. 45--46.

\section*{\S 4. Abelian varieties}

\paragraph{4.1.} Let $k$ be an algebraically closed field of characteristic $0$. Let $A$ be an abelian variety over $k$. For $n\in\mathbb N^+$, denote by $A_n$ the kernel of multiplication by $n$. The $A_n$ form a projective system
\[
\varphi_{nm,n}:A_{nm}\longrightarrow A_n,\qquad x\longmapsto x^m,
\]
and we put
\[
\widehat T(A)=\varprojlim_n A_n,\qquad
\widehat V(A)=\Af\otimes_{\widehat{\Z}}\widehat T(A).
\]
The $\widehat{\Z}$-module $\widehat T(A)$ is the product of the Tate modules, or Weil representations, $T_\ell(A)$. If $k=\C$, then
\[
\tag{4.1.1}
\widehat T(A)=\widehat{\Z}\otimes H_1(A,\Z),
\]
\[
\tag{4.1.2}
\widehat V(A)=\Af\otimes H_1(A,\Q).
\]

\paragraph{4.2.} The category of \emph{abelian varieties up to isogeny} over $k$ is the category whose objects are the abelian varieties over $k$ and in which
\[
\Hom_{\mathrm{iso}}(A,B)=\Hom(A,B)\otimes\Q.
\]
We denote by $A\otimes\Q$ the abelian variety up to isogeny underlying an abelian variety $A$. Every additive functor from the category of abelian varieties to a $\Q$-linear additive category factors through the functor $A\mapsto A\otimes\Q$. Thus the contravariant functor ``dual variety'' $A\mapsto A^\vee$ extends to abelian varieties up to isogeny. The functor $A\mapsto\widehat V(A)$ extends in the same way. Moreover, if $A_0$ is an abelian variety up to isogeny, it is equivalent to give $B$ together with an isomorphism $B\otimes\Q=A_0$, or to give $\widehat T(B)\subset\widehat V(A_0)$.

\paragraph{4.3.} Let $A$ be an abelian variety over $k$. Let $\NS(A)$ denote the group of algebraic-equivalence classes of invertible sheaves on $A$. An \emph{effective polarization}, respectively a \emph{polarization}, of $A$ is defined to be an element of $\NS(A)$, respectively of $\NS(A)\otimes\Q$, a positive multiple of which is defined by a projective embedding of $A$. A \emph{homogeneous polarization} is an element of $(\NS(A)\otimes\Q)/\Q^*$ which is the class of a polarization.

Recall that $\NS(A)\otimes\Q$ is identified, by a bijection $p\mapsto p'$, with the set of elements of $\Hom(A,A^\vee)\otimes\Q$ equal to their transpose. An isomorphism
\[
u\in\Hom(A\otimes\Q,B\otimes\Q)
\]
induces isomorphisms from $\Hom(A,A^\vee)\otimes\Q$ to $\Hom(B,B^\vee)\otimes\Q$, and from $\NS(A)\otimes\Q$ to $\NS(B)\otimes\Q$, and carries polarizations to polarizations. This makes it possible to speak of a polarization of an abelian variety up to isogeny.

\paragraph{4.4.} A polarization $p$ of $A$ defines a positive involution on $\End(A)\otimes\Q$,
\[
u\longmapsto p'^{-1}u^*p',
\]
which depends only on the homogeneous polarization $\Q^*p$.

Let $F$ be a product of totally real fields, and let $\rho:F\longrightarrow\End(A)\otimes\Q$. Let $\NS_\rho(A)\subset\NS(A)$ be the set of those $p$ such that
\[
p'\rho(f)=\rho(f)^*p'\qquad(f\in F).
\]
The vector space $\NS_\rho(A)\otimes\Q$ is endowed, for $f\in F$, with an action of $F$:
\[
(f.p)'=p'\rho(f)=\rho(f)^*p'.
\]
A \emph{weak polarization} of $A$, relative to $\rho,F$, is an element of
\[
(\NS_\rho(A)\otimes\Q)/F^*
\]
which is the class of a polarization. The set of weak polarizations of $A$ depends only on $A\otimes\Q$. The restriction to the centralizer of $F$ of the involution of $\End(A)\otimes\Q$ defined by a polarization $p\in\NS_\rho(A)$ depends only on the weak polarization $F^*p$.

\paragraph{4.5.} If $A_0^\vee$ is the dual of an abelian variety up to isogeny $A_0$, then $\widehat V(A_0)$ and $\widehat V(A_0^\vee)$ are in duality with values in $\Af(1)$. If $p$ is a polarization of $A_0$, the \emph{polarization form}
\[
\psi_p(x,y)=\langle x,p'(y)\rangle
\]
is a nondegenerate alternating form on the free $\Af$-module $\widehat V(A_0)$, with values in $\Af(1)$.

The reminders 4.1 and 4.5 extend to abelian schemes over an arbitrary base.

\paragraph{4.6.} Now take $k=\C$. The additive functor $A\mapsto H_1(A,\Q)$ extends to abelian varieties up to isogeny over $\C$ (4.2). For an abelian variety up to isogeny $A_0$, $H_1(A_0,\Q)$ and $H_1(A_0^\vee,\Q)$ are in duality. Via 4.1.2 and the isomorphism of $\Af(1)$ with $\Af$ defined by the exponential, this duality is compatible with the one considered in 4.5. The alternating form
\[
\psi_p(x,y)=\langle x,p'(y)\rangle
\]
on $H_1(A_0,\Q)$ defined by a polarization $p$ of $A_0$ is again called the polarization form.

The Hodge structure of $H_1(A_0,\Q)$ is the homomorphism
\[
h:\bS\longrightarrow \GL(H_1(A_0,\Q))_\R
\]
such that $\bS$ acts on $H_1(A_0,\Q)\otimes\C$ through the characters $z^{-1}$ and $\bar z^{-1}$, and such that $F(h)^\circ$ (1.5) is the kernel of the exponential
\[
H_1(A_0,\Q)\otimes\C\longrightarrow\Lie(A).
\]

\paragraph{Theorem 4.7 (Riemann).} The constructions
\[
(A,p)\longmapsto (H_1(A,\Q),\psi_p,H_1(A,\Z),h),
\]
\[
(V,\psi,V_\Z,h)\longmapsto
F(h)^\circ\backslash (V\otimes\C)/V_\Z
\]
establish an equivalence between
\begin{enumerate}[label=(\alph*)]
\item polarized abelian varieties over $\C$;
\item systems consisting of a vector space $V$ over $\Q$, a nondegenerate alternating form $\psi$ on $V$, an integral lattice $V_\Z$ in $V$, and a homomorphism $h$ from $\bS$ into the group $\Gp(V)_\R$ of symplectic similitudes of $V\otimes\R$, of type 1.6 and such that
\[
\psi(x,h(i)x)>0\qquad(x\in V\otimes\R,\ x\ne0).
\]
\end{enumerate}

\paragraph{Variant 4.8.} To give an abelian variety up to isogeny $A$ over $\C$, endowed with a homogeneous polarization, is equivalent to giving a vector space over $\Q$, endowed with a nondegenerate alternating form $\psi$ given up to scalar, together with $h:\bS\longrightarrow\Gp(V)$ of type 1.6.

\paragraph{4.9.} Let $L$ be a semisimple algebra with involution over $\Q$, and let $V$ be a vector space over $\Q$, endowed with a faithful $L$-module structure and with a nondegenerate alternating form $\psi$ such that
\[
\tag{4.9.1}
\psi(\ell x,y)=\psi(x,\ell^*y).
\]
Let $G$ denote the algebraic group over $\Q$ of $L$-linear symplectic similitudes of $V$. The group $G(\Q)$ is the set of $g\in\GL_L(V)$ for which there exists $\mu(g)\in\Q^*$ satisfying
\[
\tag{4.9.2}
\psi(gx,gy)=\mu(g)\psi(x,y).
\]
Suppose given a homomorphism $h_0:\bS\longrightarrow G_\R$ such that, via $h_0$, $\bS$ acts on $V_\C$ through the characters $z^{-1}$ and $\bar z^{-1}$ and such that the form $\psi(x,h_0(i)y)$ is symmetric and positive definite. The conditions (1.5.1) to (1.5.3) are then satisfied by $(G,h_0)$, and the involution of $L$ is positive.

\paragraph{4.10.} Apply (1.11) to $(G,h_0)$ and to the representation $V$ of $G$. A $G$-structure $\bar\xi$ on a vector space $H$ is interpreted, by 1.8, as the datum on $H$ of an $L$-module structure and of an alternating form $\psi$, given up to scalar, the vector space $H$, endowed with these structures, being isomorphic to $V$. Let $\bar\xi$ be a $G$-structure on $H$. A datum 1.9 (b), $h:\bS\to G^1_\R$ on $H$, then defines, by 4.8 applied to $(H,\psi,h)$, an abelian variety up to isogeny $A$, endowed with a homogeneous polarization $\bar p$, such that $H$ is $H_1(A,\Q)$, $h$ is the Hodge structure of $H_1(A,\Q)$, and $\psi$ is the polarization form. Moreover, the $L$-module structure on $H$ comes from $\rho:L\to\End(A)$, and $H$, endowed with $\bar\xi$ and $h$, is determined by $(A,\bar p,\rho)$.

For a triple $(A,\bar p,\rho)$ to arise from such an $H$, it is necessary and sufficient that
\[
\tag{4.10.1}
H_1(A,\Q),\text{ endowed with its }L\text{-module structure and with the polarization form given up to scalar, be isomorphic to }V;
\]
\[
\tag{4.10.2}
\text{for an isomorphism }t:V\longrightarrow H_1(A,\Q),\ t^{-1}h_1t\text{ be conjugate to }h_0.
\]
Let $t$ be the map from $L$ to $\C$ given by
\[
\tag{4.10.3}
t(\ell)=\Tr(\ell;V_\C/F^0(h_0)).
\]

\paragraph{Scholium 4.11.} Let $K$ be a compact open subgroup of $G(\Af)$. The points of ${}_K M_{\C}(G,h_0)$ correspond bijectively to isomorphism classes of abelian varieties up to isogeny $A$, endowed with
\begin{enumerate}[label=(\alph*)]
\item $\rho:L\to\End(A)$ such that
\[
\Tr(\rho(\ell),\Lie(A))=t(\ell)\qquad(\ell\in L),
\]
\item a homogeneous polarization $\bar p$ which induces the given involution of $L$,
\item a class modulo $K$, $\bar k$, of $L$-linear symplectic similitudes
\[
k:\widehat V(A)\xrightarrow{\sim}V\otimes\Af,
\]
\end{enumerate}
and such that
\[
(*)\qquad\text{the conditions (4.10.1) and (4.10.2) are satisfied.}
\]
One applies 1.11 and 4.10, noting that a datum 1.9 (c) is identified with a datum (c), and that the conditions in (b) and (a) result respectively from (4.10.1) and (4.10.2).

\paragraph{4.12.} Let $K$ be a compact open subgroup of $G(\Af)$, let $V_\Z$ be an integral lattice in $V$ such that $V_{\widehat{\Z}}=V_\Z\otimes\widehat{\Z}$ is $K$-invariant, let $L_\Z$ be an order in $L$ such that $L_\Z V_\Z\subset V_\Z$, let $\psi_\Z$ be a positive rational multiple of the alternating form of $V$, with integral values on $V_\Z$, and let $V'_\Z$ be the largest lattice in $V$ such that $\psi_\Z(V_\Z,V'_\Z)\subset \Z(1)$. Let $n$ be an integer, and put
\[
K_n=\{\,g\in G(\Af)\mid (g-1)V'_{\widehat{\Z}}\subset nV_{\widehat{\Z}}\,\}.
\]
Assume that $n$ is large enough that $K\supset K_n$.

If $A$ is as in 4.11, then $k^{-1}(V_{\widehat{\Z}})\subset\widehat V(A)$ does not depend on $k\in\bar k$, and defines an abelian variety $B$ with $B\otimes\Q\simeq A$ and
\[
\widehat T(B)=k^{-1}(V_{\widehat{\Z}})\subset\widehat V(A).
\]
The action of $L$ on $A$ induces an action of $L_\Z$ on $B$. There exists a unique effective polarization $p$ of $B$, with $p\in\bar p$, such that $k^{-1}(V'_{\widehat{\Z}})$ is the largest ``lattice'' $\widehat T'(B)\supset\widehat T(B)$ satisfying
\[
\psi_p(\widehat T(B),\widehat T'(B))\subset \widehat{\Z}(1).
\]
Finally, the set $\bar k$ of $L$-linear symplectic isomorphisms from $\widehat T(B)$ to $V_{\widehat{\Z}}$ is the inverse image of its image $\bar k_n$ in $\Isom(B_n,V_\Z/nV_\Z)$.

This again permits the points of ${}_K M_{\C}(G,h_0)$ to be interpreted as corresponding to isomorphism classes of systems $(B,p,\rho,\bar k_n)$ consisting of
\begin{enumerate}[label=(\alph*)]
\item a polarized abelian variety $(B,p)$, with complex multiplication by $L_\Z$, satisfying 4.11 (a) and (b);
\item a class modulo $K/K_n$, $\bar k_n$, of isomorphisms
\[
k_n:B_n\xrightarrow{\sim}V_\Z/nV_\Z,
\]
which can be lifted to $L$-linear symplectic isomorphisms
\[
k:\widehat T(B)\longrightarrow V_{\widehat{\Z}},
\]
and which satisfy (4.10.1) and (4.10.2).
\end{enumerate}

\paragraph{4.13.} Let $F$ denote the algebra of invariants under $*$ in the center of $L$. It is a product of totally real fields. There exists a unique $F$-bilinear alternating form $\psi_F$ on $V$ such that
\[
\Tr_{F/\Q}\psi_F=\psi.
\]
Let $G_1$ denote the group of $L$-linear symplectic similitudes of $\psi_F$: $G_1(\Q)$ is the set of $g\in\GL_L(V)$ for which there exists $\mu(g)\in F^*$ satisfying
\[
\tag{4.13.1}
\psi_F(gx,gy)=\mu(g)\psi_F(x,y),
\]
that is,
\[
\tag{4.13.2}
\psi(gx,gy)=\psi(\mu(g)x,y).
\]
As above, one verifies the following variant.

\paragraph{Variant 4.14.} Let $K$ be a compact open subgroup of $G_1(\Af)$. The points of ${}_K M_{\C}(G_1,h_0)$ correspond bijectively to isomorphism classes of abelian varieties up to isogeny $A$, endowed with
\begin{enumerate}[label=(\alph*)]
\item $\rho:L\longrightarrow\End(A)$ as in 4.11,
\item a weak polarization, relative to $F$, $\bar p$, which induces the given involution of $L$,
\item a class modulo $K$ of $L$-linear symplectic similitudes
\[
k:\widehat V(A)\xrightarrow{\sim}V\otimes\Af
\]
for the $F\otimes\Af$-bilinear forms on the two sides,
\end{enumerate}
and such that
\[
(*)\qquad\text{conditions analogous to (4.10.1) and (4.10.2) are satisfied.}
\]

\paragraph{4.15.} One can make 4.11 and 4.14 more precise by showing that ${}_K M_{\C}(G,h)$ is the coarse moduli scheme of the evident functor on schemes of finite type over $\C$.

\paragraph{4.16. Example I} (continuation of 1.6, 1.11). Place ourselves in the special case of 4.9 where $L=\Q$ and $\dim(V)=2g$. Then $G=\Gp(V)$, the group of symplectic similitudes. For $h_0:\bS\to G_\R$ of type 1.6, one has $E(\Gp,h_0)=\Q$. Let $V_\Z$, $n$, and $K_n$ be as in 1.11. For every scheme $S$, let $F(S)$ be the set of isomorphism classes of principally polarized abelian schemes $A/S$, endowed with a symplectic similitude
\[
k_n:A_n\xrightarrow{\sim}(V_\Z/nV_\Z)_S.
\]
For $n\ge3$, the classified objects no longer have automorphisms and the functor $F$ is represented by a $\Q$-scheme ${}_{K_n}M$ [6]. From 1.11, 4.11, and 4.12, one has
\[
m_n:{}_{K_n}M(\C)\xrightarrow{\sim}{}_{K_n}M_{\C}(\Gp,h_0).
\]
In the language of 3.1, this gives:

\paragraph{Proposition 4.17.} The $\Q$-scheme
\[
M(\Gp,h_0)=\varprojlim_{K_n}{}_{K_n}M
\]
is a model over $\Q=E(\Gp,h_0)$ of $M_{\C}(\Gp,h_0)$.

\paragraph{4.18.} Place ourselves in the special case of 4.9 where $V$ is a monogenic $L$-module. The groups $G$ and $G_1$ are then contained in the center of $L^*$.

Let $\bar E$ be an algebraic closure of $E(G,h)$ and let $M^+$ be the set of isomorphism classes $[A,\rho,k,\bar p]$ of objects $(A,\rho,k,\bar p)$ consisting of
\begin{enumerate}[label=(\alph*)]
\item an abelian variety up to isogeny $A/\bar E$, endowed with a homogeneous polarization $\bar p$ and with $\rho:L\to\End(A)$ such that, with the notation of 4.10.3, for $\ell\in L$ one has
\[
\Tr(\rho(\ell),\Lie(A))=t(\ell)\in E(G,h);
\]
\item an $L\otimes\Af$-linear isomorphism
\[
k:\widehat V(A)\xrightarrow{\sim}V\otimes\Af.
\]
\end{enumerate}
The abelian variety $A$ is therefore of CM type. We refer to [16] for the proof of the following fundamental theorem.

\paragraph{Theorem 4.19 (Shimura--Taniyama).} The Galois group $\Gal(\bar E/E(G,h))$ acts on $M^+$ through its largest abelian quotient
\[
\pi_0\bigl(E(G,h)^*(\A)/E(G,h)^*(\Q)\bigr).
\]
For $e\in E(G,h)^*(\A)$, with image $\varphi(e)$ in abelianized Galois and with component $e^f$ in $E(G,h)^*(\Af)$, one has
\[
\tag{4.19.1}
\varphi(e)([A,\rho,k,\bar p])=[A,\rho,r(G,h)(e^f)k,\bar p].
\]

Let $\tau$ be an embedding of $\bar E$ into $\C$ extending that of $E(G,h)$, let $K$ be a compact open subgroup of $G(\Af)$, and let ${}_K M(G,h)(\bar E)$ be the set of isomorphism classes of objects $(A,\rho,\bar k,\bar p)$ where
\begin{enumerate}[label=(\alph*)]
\item $A,\rho$, and $\bar p$ are as above;
\item $\bar k$ is a class modulo $K$ of $L\otimes\Af$-linear symplectic similitudes
\[
k:\widehat V(A)\xrightarrow{\sim}V\otimes\Af;
\]
\item the vector space $H_1(A_\C,\Q)$, where $A_\C$ is defined via $\tau$, endowed with the action of $L$ and with the polarization form given up to scalar, is isomorphic to $V$.
\end{enumerate}
By 4.19, condition (c) is independent of the choice of $\tau$, so that the finite set ${}_K M(G,h)(\bar E)$ is endowed with an action of $\Gal(\bar E/E(G,h))$; as a Galois set, it is identified with the set of $\bar E$-points of a finite étale $E(G,h)$-scheme ${}_K M(G,h)$. By 4.14, ${}_K M(G,h)(\C)=G(\R)XK\backslash G(\A)/G(\Q)$, and by 4.19 one has:

\paragraph{Proposition 4.20.} The $E(G,h)$-scheme
\[
M(G,h)=\varprojlim_K {}_K M(G,h)
\]
is a canonical model of $M_{\C}(G,h)$.

\paragraph{Theorem 4.21.} The model $M(\Gp,h_0)$ constructed in 4.17 is a canonical model.

Let $(V,\psi)$ be as in 4.16, and let $L$, $\rho:L\to\End(V)$, $G$, and $h:\bS\to G_\R$ be as in 4.18. We have
\[
u:(G,h)\longrightarrow(\Gp,h_0).
\]
Let $K\subset G(\Af)$ be such that $u(K)\subset K_n$ (4.16). If $\bar E$ is an algebraic closure of $E(G,h)$, define
\[
u:{}_K M(G,h)(\bar E)\longrightarrow{}_{K_n}M(\Gp,h_0)(\bar E)
\]
by ``forgetting the complex multiplication'', associating to $[A,\rho,\bar k,\bar p]$ the triple consisting of
\begin{enumerate}[label=(\alph*)]
\item the abelian variety $B$, endowed with $B\otimes\Q\simeq A$ and such that
\[
\widehat T(B)=k^{-1}(V_{\widehat{\Z}})\subset\widehat V(A)\qquad(k\in\bar k);
\]
\item the polarization $\bar p$ of $B$;
\item the isomorphism
\[
k:B_n\xrightarrow{\sim}V_\Z/nV_\Z\qquad(k\in\bar k).
\]
\end{enumerate}
The maps $u$ define a morphism of models
\[
u:M(G,h)\longrightarrow M(\Gp,h_0)\otimes E(G,h).
\]
The proof is completed by showing ([8]) that for every
\[
u:(H,h')\longrightarrow(\Gp,h_0)
\]
as in 3.13, there exists $v:(G,h)\longrightarrow(\Gp,h_0)$ as above, for a suitable $L$, such that $uh'=vh$.

\paragraph{Variant 4.22.} Let $F$ be a totally real number field, let $n=[F:\Q]$, and let $\Gp_{2g}^{(F)}$ be the group of symplectic similitudes of a symplectic $F$-vector space of dimension $2g$. Let
\[
h_0:\bS\longrightarrow \Gp_{2g,\R}^{(F)}\simeq \Gp_{2g}(\R)^n
\]
be as in 1.6. One can construct a canonical model of $M_{\C}(\Gp_{2g}^{(F)},h_0)$ from moduli of weakly polarized abelian schemes.





\section*{\S 5. Construction techniques}

Definition 3.13 is usable only because there are ``many'' special points.

\paragraph{Theorem 5.1.} Let $G$ and $h$ be as in 3.13. For every finite extension $F$ of $E(G,h)$, there exists $(H,h',u:H\hookrightarrow G)$ as in 3.13 such that the extension $E(H,h')$ of $E(G,h)$ is linearly disjoint from $F$.

Let
\[
Y=\Hom(\Gm,G)/G
\]
be the scheme of conjugacy classes of morphisms from $\Gm$ to $G$. The Galois group $\Gal(\overline{\Q}/\Q)$ acts on the infinite discrete set $Y(\C)$ of conjugacy classes of morphisms from $\Gm{}_{\C}$ to $G_{\C}$. Let $Y_o$ be the finite subscheme of $Y$ such that $Y_o(\C)$ is the Galois orbit of the class of $hr$. It is the spectrum of $E(G,h)$.

Let $V'=\Lie(G)$, regarded as an affine space over $\Q$, and let $V$ be the open subscheme of regular elements of the Lie algebra. For $v$ in $V$, let $T_v$ denote the centralizer of $v$, a maximal torus. Let $W$ be the moduli scheme of maximal tori $T$ of $G$, endowed with
\[
s:\Gm\longrightarrow T,\qquad v\in\Lie(T),
\]
such that
\begin{enumerate}[label=(\alph*)]
\item $v$ is regular in $\Lie(G)$;
\item the class of $s:\Gm\to G$ lies in $Y_o$.
\end{enumerate}
The morphism
\[
f:W\longrightarrow V,\qquad (T,s,v)\longmapsto v
\]
is finite étale and surjective; note that $T=T_v$. Let $p$ be the morphism from $W$ to $Y_o$,
\[
p:(T,s,v)\longmapsto \text{the class of }s.
\]
Thus one has the diagram
\[
\tag{5.1.1}
\begin{array}{ccc}
W & \xrightarrow{\ f\ } & V\\
{\scriptstyle p}\downarrow && \downarrow\\
\Spec(E(G,h)) & \longrightarrow & \Spec(\Q).
\end{array}
\]

\paragraph{Lemma 5.1.2.} The scheme $W$ is irreducible, and the geometric fibers of $p$ are geometrically irreducible.

Since $Y_o$ is irreducible, being the spectrum of an extension of $\Q$, it is enough to prove the second assertion. It follows from the following facts:
\begin{enumerate}[label=(\alph*)]
\item since $G$ is connected, the scheme over $\C$ of homomorphisms from $\Gm{}_{\C}$ to $G_{\C}$ conjugate to a given homomorphism is irreducible;
\item for a given $i:\Gm{}_{\C}\to G_{\C}$, by the theorem on conjugacy of maximal tori applied to the connected centralizer of $i(\Gm{}_{\C})$, the scheme of maximal tori of $G_{\C}$ containing $i(\Gm{}_{\C})$ is irreducible.
\end{enumerate}

Let $U\subset V(\R)$ be the set of $v\in V(\R)$ such that $(T_v/C)(\R)$ is compact. For the usual topology, $U$ is open.

For $v\in U$ and $w\in W(\C)$ such that $f(w)=v\in V(\R)\subset V(\C)$, there exists a unique homomorphism
\[
h(w):\bS\longrightarrow T_v
\]
such that $h(w)\circ r=s_w$: after extension of scalars to $\C$,
\[
h(w)=s_w\circ z\cdot \overline{s_w}\circ\overline z.
\]
Let $v\in U$. By the theorem on conjugacy of maximal compact subgroups in $G/C(\R)^\circ$, there exists $g\in G(\R)$, even $g\in G'(\R)^\circ$, such that
\[
gT_vg^{-1}\subset K_\infty.
\]
Since $T_v$ is a maximal torus, $gT_vg^{-1}$ contains the center of $K_\infty$. It follows that there exists $w\in W(\C)$ such that $f(w)=v$ and such that $h(w)$ is conjugate to $h$.

The proof of 5.1 is concluded by applying to the diagram (5.1.1) and to $U$ the following variant of Hilbert's irreducibility theorem, for whose proof I am indebted to J.-P. Serre.

\paragraph{Lemma 5.13.} Let $E$ be a finite extension of $\Q$, and let
\[
\begin{array}{ccc}
W & \xrightarrow{\ f\ } & V\\
{\scriptstyle p}\downarrow && \downarrow\\
\Spec(E) & \longrightarrow & \Spec(\Q)
\end{array}
\]
be a commutative diagram in which
\begin{enumerate}[label=(\roman*)]
\item $V$ is a Zariski open subset of an affine space over $\Q$;
\item the geometric fibers of $p$ are irreducible, and $W$ is reduced;
\item $f$ is quasi-finite and dominant.
\end{enumerate}
Then, for every open set $U\subset V(\R)$ in the usual topology and every extension $F$ of $E$, there exists
\[
v\in V(\Q)\cap U
\]
such that $f^{-1}(v)$ is the spectrum of an extension of $E$ linearly disjoint from $F$.

Let $W'=W\otimes_EF$. The hypotheses of 4.12 are again satisfied for $W'/F$ and
\[
f':W'\longrightarrow W\longrightarrow V,
\]
and one must construct $v\in V(\Q)\cap U$ such that
\[
f'^{-1}(v)=f^{-1}(v)\otimes_EF
\]
is the spectrum of a field. This reduces us to the case $E=F$, which follows from Hilbert's irreducibility theorem; cf. Lang, \emph{Diophantine Geometry}, chap. VIII.

\paragraph{Proposition 5.2.} For $x\in\MC(G,h)$ and $K$ compact open in $G(\Af)$, the image in ${}_K\MC(G,h)$ of $G(\Af)\cdot x\subset\MC(G,h)$ is dense.

This is a consequence of (0.4). The reader may verify it.

\paragraph{Proposition 5.3.} Let $E$ be a field endowed with a complex embedding $\rho$, let $X$ be a scheme over $E$, and let $Y$ be a reduced closed subscheme of $X_{\C}$. Suppose that there exist finite extensions $E_i$ of $E$ in $\C$, schemes $M_{i,\alpha}/E_i$, and $E_i$-morphisms
\[
v_{i,\alpha}:M_{i,\alpha}\longrightarrow X\otimes_EE_i,
\]
or, writing
\[
v_i:\coprod_\alpha M_{i,\alpha}\longrightarrow X\otimes_EE_i,
\]
such that $E=\bigcap_i E_i$,
\[
v_i\Bigl(\bigl(\coprod_\alpha M_{i,\alpha}\bigr)\otimes_{E_i}\C\Bigr)\subset Y,
\]
and
\[
v_i\Bigl(\bigl(\coprod_\alpha M_{i,\alpha}\bigr)\otimes_{E_i}\C\Bigr)
\]
is dense in $Y$. Then $Y$ is defined over $E$.

\paragraph{Corollary 5.4.} Let $u:(G^1,h^1)\to(G^2,h^2)$ be as in 1.14. Let $M_E(G^1,h^1)$ and $M_E(G^2,h^2)$ be weakly canonical models of the $\MC(G^i,h^i)$ over a number field $E$, with $E(G^i,h^i)\subset E\subset\C$. Then
\[
v:\MC(G^1,h^1)\longrightarrow\MC(G^2,h^2)
\]
is defined over $E$, cf. 3.12.

Let $K^i\subset G^i(\Af)$ be compact open subgroups such that $v(K^1)\subset K^2$. We prove that
\[
v(K^1,K^2):{}_{K^1}\MC(G^1,h^1)\longrightarrow{}_{K^2}\MC(G^2,h^2)
\]
is defined over $E$. Let $u:(H,h')\to(G^1,h^1)$ be as in 3.13 and set $F=E\cdot E(H,h')$. For $g\in G^1(\Af)$, one has $F$-morphisms
\[
\begin{aligned}
a(g):\quad M_F(H,h')&\longrightarrow M_F(G^1,h^1)
 \xrightarrow{\ g\ } M_F(G^1,h^1)
 \longrightarrow {}_{K^1}M_F(G^1,h^1),\\
b(g):\quad M_F(H,h')&\longrightarrow M_F(G^2,h^2)
 \xrightarrow{\ v(g)\ } M_F(G^2,h^2)
 \longrightarrow {}_{K^2}M_F(G^2,h^2),
\end{aligned}
\]
and, after extension of scalars to $\C$,
\[
v(K^1,K^2)\circ a(g)=b(g).
\]
By 5.1 and 5.2, one can apply 5.3 to
\[
X={}_{K^1}M_E(G^1,h^1)\times{}_{K^2}M_E(G^2,h^2),
\]
to the graph $Y$ of $v$, to the $E_i=E\cdot E(H,h')$ for variable $(H,h',u)$, and to
\[
(a(g),b(g)):M_{E_i}(H,h')\longrightarrow X_{E_i}.
\]
It follows that the graph of $v$ is defined over $E$.

In the special case $(G^1,h^1)=(G^2,h^2)$ and $u=\Id$, 5.4 says:

\paragraph{Corollary 5.5.} Up to unique isomorphism, there exists at most one weakly canonical model of $\MC(G,h)$ over $E$, for $E(G,h)\subset E\subset\C$.

\paragraph{Corollary 5.6.} If $G'$ is simply connected, the reciprocity law of a weakly canonical model of $\MC(G,h)$ over $E$, for $E(G,h)\subset E\subset\C$, is given by 3.9.1.

This is the same proof as in the special case $v:(G,h)\to(T,vh)$ of 5.4.

\paragraph{Corollary 5.7.} Let $v:(G^1,h^1)\hookrightarrow(G^2,h^2)$ be as in 1.15. If $\MC(G^2,h^2)$ admits a weakly canonical model over $E$, with
\[
E(G^2,h^2)\subset E(G^1,h^1)\subset E\subset\C,
\]
then $\MC(G^1,h^1)$ admits a weakly canonical model over $E$.

Let $K^i$ be a compact open subgroup of $G^i(\Af)$. Using 1.15, we are reduced to showing only that, if $v(K^1)\subset K^2$ and
\[
{}_{K^1}\MC(G^1,h^1)\subset{}_{K^2}\MC(G^2,h^2),
\]
then the subscheme ${}_{K^1}\MC(G^1,h^1)$ of
\[
{}_{K^2}\MC(G^2,h^2)={}_{K^2}M_E(G^2,h^2)\otimes_E\C
\]
is defined over $E$.

Let $u:(H,h')\to(G^1,h^1)$ be as in 3.13 and, for $g\in G^1(\Af)$, let $a(g)$ be the composite homomorphism, defined over $F=E\cdot E(H,h')$:
\[
a(g):M_F(H,h')\longrightarrow M_F(G^1,h^1)
 \xrightarrow{\ v(g)\ } M_F(G^2,h^2)
 \longrightarrow{}_{K^2}M_F(G^2,h^2).
\]
By 5.1 and 5.2, one concludes by applying 5.3 to
\[
X={}_{K^2}M_E(G^2,h^2),\qquad
Y={}_{K^1}\MC(G^1,h^1),
\]
to the $E_i=E\cdot E(H,h')$ for variable $(H,h',u)$ and to the
\[
a(g):M_{E_i}(H,h')\longrightarrow X_{E_i}.
\]

\paragraph{Example 5.8.} For $(G,h_0)$ as in 4.9, one constructs, by 5.7 and 4.21, a canonical model
\[
M(G^\circ,h_0)\subset M(\Gp(V,\psi),h_0)\otimes_{\Q}E(G^\circ,h^\circ).
\]

\paragraph{Variant 5.9.} Let $(G_1,h_0)$ be such that, with the notation of 4.22, there exists
\[
u:(G_1,h_0)\hookrightarrow(\Gp_{2g}^{(F)},h_0),
\]
for example $(G_1^\circ,h_0)$ with $(G_1,h_0)$ as in 4.13. An application of 5.7 gives a canonical model of $M(G_1,h_0)$.

\paragraph{Proposition 5.10.} Let $(G,h)$ be as in 3.13, let
\[
E_i:\ E(G,h)\subset E_i\subset\C
\]
be number fields, let $E$ be the intersection of the $E_i$, and let $M_{E_i}(G,h)$ be a weakly canonical model of $\MC(G,h)$ over $E_i$. Then there exists a unique model $M_E(G,h)$ of $\MC(G,h)$ over $E$ such that, for every $i$, the model $M_E(G,h)\otimes_EE_i$ over $E_i$ is isomorphic to $M_{E_i}(G,h)$.

In view of the uniqueness assertion, one may suppose the $E_i$ finite in number. Let $F$ be the extension of $E$ in $\C$ that they generate, and let $\mathfrak E$ be the set of subextensions of $F$ containing one of the $E_i$. Using 5.5, one obtains for every $F'\in\mathfrak E$ a model $M_{F'}(G,h)$, and for every inclusion $F'\subset F''$, $F',F''\in\mathfrak E$, an isomorphism of models
\[
\alpha(F'',F'):M_{F'}(G,h)\otimes_{F'}F''\longrightarrow M_{F''}(G,h).
\]

Let $u:(H,h')\to(G,h)$ be as in 3.13, with $E(H,h')$ linearly disjoint from $F$ over $E(G,h)$, cf. 5.1. Put
\[
E_1=E(H,h')\cdot E,\qquad F'_1=E_1\cdot F'=E(H,h')\cdot F'\qquad(F'\in\mathfrak E).
\]
Let $g\in G(\Af)$ and let $K\subset G(\Af)$ be a compact open subgroup. Since $M_{F'}(G,h)$ is weakly canonical, one has an $F'_1$-morphism
\[
a'(g):M_{F'_1}(H,h')\longrightarrow M_{F'_1}(G,h)
 \xrightarrow{\ g\ } M_{F'_1}(G,h)
 \longrightarrow {}_K M_{F'_1}(G,h).
\]
Since $F'_1=E_1\otimes_EF'$, one may ``interpret'' this as a morphism of $F'$-schemes
\[
a(g):M_{E_1}(H,h')\otimes_EF'\longrightarrow {}_K M_{F'}(G,h).
\]
The $a(g)$ satisfy
\[
\alpha(F'',F')\,a(g)_{F''}=a(g).
\]
One concludes by 5.2 and by the following lemma, applied to
\[
X^1=\coprod_{g\in G(\Af)}M_{E_1}(H,h').
\]
Its proof is left to the reader as an exercise.

\paragraph{Lemma 5.10.1.} Let $F/E$ be a finite extension of fields and let $\mathfrak E$ be a set of subfields of $F$, whose intersection is reduced to $E$, such that
\[
F\in\mathfrak E,\qquad F'\subset F''\subset F\Longrightarrow F''\in\mathfrak E.
\]
\begin{enumerate}[label=(\roman*)]
\item The functor which sends a scheme $X/E$ to the system $\mathcal X$ of schemes $X_{F'}/F'$, for $F'\in\mathfrak E$, and of isomorphisms
\[
X_{F'}\otimes_{F'}F''\xrightarrow{\sim}X_{F''}\qquad(F'\subset F'',\ F',F''\in\mathfrak E)
\]
is fully faithful.
\item For a system
\[
\mathcal X=\{X(F')/F'\ (F'\in\mathfrak E);\ \alpha(F'',F'):X(F')\otimes_{F'}F''\xrightarrow{\sim}X(F'')\}
\]
to come from a scheme $X/E$, it is enough that the $X(F')/F'$ be quasi-projective and that there exist a scheme $X^1/E$ and a morphism
\[
u:X^1_{\mathfrak E}\longrightarrow\mathcal X
\]
such that the $u(F'):X^1_{F'}\to X(F')$ have dense image.
\end{enumerate}

\paragraph{Remark 5.10.2.} Under the hypotheses of 5.10, for $u:(H,h')\hookrightarrow(G,h)$ as in 3.13, the field of definition of
\[
u:\MC(H,h')\longrightarrow\MC(G,h)
\]
is contained in the extension
\[
\bigcap_i E(H,h')\cdot E_i
\]
of $E(H,h')\cdot E$.

\paragraph{Proposition 5.11.} Let $h:\bS\to G_\R$ be as in 3.13 and let $\delta:\bS\to C^\circ_\R$, where $C$ is the center of $G$. If
\[
E(G,h),\ E(C,\delta)\subset E\subset\C
\]
and if $\MC(G,h)$ admits a weakly canonical model over $E$, then $\MC(G,h\delta)$ also admits one.

Let $K$ be a compact open subgroup of $G(\Af)$ and let $L=K\cap C^\circ(\Af)$. Over $\C$, the morphism
\[
u:(G\times C^\circ,(h,\delta))\longrightarrow(G,h\delta),\qquad (g,c)\longmapsto gc,
\]
induces morphisms
\[
u_K:{}_K\MC(G,h)\times {}_L\MC(C^\circ,\delta)
 \longrightarrow {}_K\MC(G,h\delta).
\]
Let $d$ be the morphism
\[
d:C^\circ\longrightarrow G\times C^\circ,\qquad c\longmapsto(c,c^{-1}).
\]
The subgroup $d(C^\circ(\Af))$ of $G\times C^\circ(\Af)$ acts on
\[
\MC(G,h)\times\MC(C^\circ,\delta)
\]
through the quotient $\pi_0(C^\circ(\A))$ of $C^\circ(\Af)$, and $u_K$ induces an isomorphism
\[
\bar u_K:\bigl({}_K\MC(G,h)\times{}_L\MC(C^\circ,\delta)\bigr)/
\bigl(L\backslash\pi_0(C^\circ(\A))\bigr)
\xrightarrow{\sim}{}_K\MC(G,h\delta).
\]
Define ${}_K M_E(G,h\delta)$ to be the quotient
\[
\bigl({}_K M_E(G,h)\times{}_L M_E(C^\circ,\delta)\bigr)/
\bigl(L\backslash\pi_0(C^\circ(\A))\bigr).
\]
One verifies that this gives a weakly canonical model. By construction, one has
\[
\tag{5.11.1}
M_E(G,h)\times_E M_E(C^\circ,\delta)\longrightarrow M_E(G,h\delta).
\]

\paragraph{Conjecture 5.12.} For every $(G,h_0)$ as in 3.13, there exists a canonical model of $\MC(G,h_0)$.

Suppose, this being a typical case, that $G$ is $\Q$-simple and adjoint. Let $\widetilde G$ be the universal covering of $G$ and let $\omega$ be a weight of $\widetilde G_{\C}$. If, with the notation of 3.7, the conjugacy class of $h_0r$ is described by integers $(n_i)_{i\in|\Delta|}$, and if
\[
\omega=\sum m_i\omega_i,
\]
put
\[
\langle\omega,h_0r\rangle=\sum n_i m_i.
\]
Up to now one has been able to attack 5.12 only when $\widetilde G$ admits a nontrivial representation $V$ such that all irreducible components $V'$ of $V_{\C}$ satisfy:
\[
(*)\qquad
\text{as }\omega\text{ runs through the weights of }V',\
\langle\omega,h_0r\rangle\text{ takes at most two values.}
\]
Let $\sigma$ be the opposition involution. Condition $(*)$ also means that the dominant weight $\omega$ of $V'$ satisfies
\[
\langle\omega,h_0r\rangle+\langle\sigma\omega,h_0r\rangle=0\text{ or }1.
\]
There is no such representation $V$ if $G$ is exceptional, in which case the noncompact factors of $G_{\R}$ are of type $E_{6(-14)}$ or $E_{7(-25)}$, or of trialitarian type $D_4$, nor if $G$ is of type $D_n$ and $G_{\R}$ has noncompact factors of both types $D_n^{\mathbb H}$ and $D_n^{\mathbb R,2}$.

\section*{\S 6. Strange models}

\paragraph{6.1.} Let $F$ be a totally real number field, and let
\[
\tag{6.1.1}
{}_F\Sigma:\quad 0\longrightarrow {}_FG'\longrightarrow {}_FG
 \longrightarrow \Gm\longrightarrow0
\]
be a form over $F$ of the exact sequence
\[
\tag{6.1.2}
{}_F\Sigma_0:\quad 0\longrightarrow\Sp_{2n}\longrightarrow\Gp_{2n}
 \xrightarrow{\ \mu\ }\Gm\longrightarrow0,
\]
where $\Gp$ denotes a group of symplectic similitudes.

Assume that, for every real place $\tau$ of $F$, $\Sigma\otimes_{F,\tau}\R$ is of one of the following types:
\begin{enumerate}[label=(\alph*)]
\item ${}_F\Sigma\otimes_{F,\tau}\R$ is isomorphic to ${}_F\Sigma_0\otimes_{F,\tau}\R$, or
\item $G'\otimes_{F,\tau}\R$ is compact.
\end{enumerate}
Let $\tau_1,\ldots,\tau_g$ denote the real places of $F$, and suppose that the places of type (a) are $\tau_1,\ldots,\tau_r$, with $r>0$.

\paragraph{6.2.} The group ${}_FG$ may be described as follows:
\begin{enumerate}[label=(\alph*)]
\item Choose a quaternion algebra $B$ over $F$, indefinite at the places $\tau_i$ for $i\le r$, and definite for $i>r$. Let $x\mapsto\overline{x}$ denote its canonical involution.
\item Choose a free $B$-module $V$, endowed with a nondegenerate $F$-bilinear form $\Phi$, symmetric and such that
\[
\Phi(bx,y)=\Phi(x,\overline b\,y).
\]
\item Assume that, for $i>r$, the form deduced from $\Phi$ by extension of scalars through $\tau_i:F\to\R$ is positive definite.
\item Take ${}_FG$ to be the group of $B$-linear similitudes of $\Phi$; in other words, the elements $g\in{}_FG(F)$ are the $g\in\GL_B(V)$ for which there exists $\mu(g)\in F^*$ satisfying
\[
\Phi(gx,gy)=\mu(g)\Phi(x,y).
\]
\end{enumerate}

\paragraph{6.3.} Let
\[
\Sigma:\quad 0\longrightarrow G'\longrightarrow G\longrightarrow F^*\longrightarrow0
\]
denote the exact sequence of algebraic groups over $\Q$ deduced from ${}_F\Sigma$ by Weil restriction of scalars from $F$ to $\Q$. We have
\[
\tag{6.3.1}
G_{\R}=\prod_{\tau}{}_FG\otimes_{F,\tau}\R
\qquad(\tau\text{ a real place}).
\]
A morphism $h:\bS\to G_\R$ is defined by its coordinates
\[
h_i:\bS\longrightarrow {}_FG\otimes_{F,\tau_i}\R.
\]
By 1.6, there exists a unique conjugacy class of morphisms $h$ such that, for $i\le r$, $h_i$ is of type 1.6, and for $i>r$, $h_i$ is trivial. Let $h_0$ be in this class.

The field $E(G,h_0)$ is the totally real subfield of $\C$ generated by the elements
\[
\sum_{i=1}^r\tau_i(f)\qquad(f\in F).
\]

\paragraph{Theorem 6.4.} With the preceding notation, $\MC(G,h_0)$ admits a canonical model $M(G,h_0)$.

Let $Z$ be a totally imaginary quadratic extension of $F$. Let
\[
V'=V\otimes_F Z,\qquad L=B\otimes_F Z,
\]
and let $G'$ be the algebraic subgroup of $\GL_L(V')$ generated by $G$ and $Z^*$:
\[
\tag{6.4.1}
0\longrightarrow F^*\xrightarrow{\ (f,f^{-1})\ }G\times Z^*
 \xrightarrow{\ q\ }G'\longrightarrow0.
\]
Let $x\mapsto\widetilde x$ be the involution of $L$ given by the tensor product of the involutions of $B$ and $Z$. Up to a factor in $F$, there exists on $Z$ a unique alternating $F$-bilinear form $\psi_Z$. If ${}_F\psi$ is the alternating $F$-bilinear form on $V'$ which is the tensor product of $\Phi$ and $\psi_Z$, then
\[
{}_F\psi(\ell x,y)={}_F\psi(x,\widetilde\ell\,y).
\]
Put
\[
\psi(x,y)=\Tr_{F/\Q}\bigl({}_F\psi(x,y)\bigr).
\]

Choose, for every real place of $F$, a complex embedding of $Z$ inducing it. Then
\[
Z^*(\R)\simeq\prod_{i=1}^g \C^*.
\]
Let $h_Z$ be the homomorphism from $\bS$ into $Z^*_\R$ with coordinates
\[
h_i:\bS(\R)=\C^*\longrightarrow\C^*\qquad(1\le i\le g)
\]
equal to $z\mapsto1$ for $i\le r$, and to $z\mapsto z^{-1}$ for $i>r$.

For every $h:\bS\to G_\R$ as in 6.3, put
\[
h'=q\circ(h\times h_Z):\bS\longrightarrow G'_\R.
\]

\paragraph{Lemma 6.4.2.} There exists $b\in L$ such that, for a suitable choice of $\psi_Z$, the form
\[
\psi\bigl(x,b\,h'_0(i)y\bigr)\qquad\text{on }V'\otimes_\Q\R
\]
is symmetric and positive definite.

This is verified after extension of scalars from $\Q$ to $\R$.

This lemma and 5.9 make it possible to construct a canonical model of $M(G',h'_0)$, since $G'$ is contained in the group of symplectic similitudes of ${}_F\psi(x,by)$. Applying 5.11 and 5.7, one obtains a weakly canonical model of $\MC(G,h_0)$ over the extension $E(Z^*,h_Z)$ of $E(G,h_0)$. The conclusion follows from 5.10, 5.10.2 and the following lemma, which can be deduced from 5.13.

\paragraph{Lemma 6.5.} Let $F$ be a totally real number field, endowed with a set $S$ of real embeddings. Let $F^\ast$ be the number field corresponding to the subgroup of $\Gal(\overline{\Q}/\Q)$ which stabilizes $S$. For $Z$ a totally imaginary quadratic extension of $F$, endowed with a set $T$ of complex embeddings, one above each element of $S$, let likewise $Z^\ast\supset F^\ast$ be the field of invariants of the stabilizer of $T$. Then, for every extension $E$ of $F^\ast$, there exist $Z$ and $T$ such that $Z^\ast$ is linearly disjoint from $E$.

\paragraph{Remark 6.6.} The canonical map 1.13, 1.14
\[
\tag{6.6.1}
\MC(G,h_0)\times\MC(Z^*,h_Z)\longrightarrow\MC(G',h'_0)
\]
is interpreted as follows. Each conjugate of $h$ defines on $V_{\C}$ a




a Hodge bigrading invariant under $F$:
\[
V_{\C}=V^{-1,0}\oplus V^{0,-1}\oplus V^{0,0}
\qquad(\overline{V^{pq}}=V^{qp}).
\]
One should beware that the grading, by weight, of $V_{\C}$ by
\[
V^n=\sum_{p+q=n}V^{pq}
\]
is not defined over $\Q$ if $r\ne g$.

Similarly, $Z_{\C}$ is bigraded:
\[
Z_{\C}=Z^{-1,0}\oplus Z^{0,-1}\oplus Z^{0,0}.
\]
The miracle is that the $\Q$-vector space $V'=V\otimes_F Z$, endowed with the tensor-product Hodge bigrading
\[
V'_{\C}=V_{\C}^{\prime -1,0}\oplus V_{\C}^{\prime 0,-1}\qquad(!)
\]
is the $H_1$ of an abelian variety.

\section*{BIBLIOGRAPHY}

\begin{description}[leftmargin=3.0em,labelwidth=2.4em,style=nextline]
\item[{[1]}] F. Bruhat and J. Tits -- \emph{Groupes algébriques simples sur un corps local}, Proc. on a conf. on local fields, p. 23--26, Driebergen 1966, Springer-Verlag.
\item[{[2]}] P. Deligne -- \emph{Travaux de Griffiths}, Sém. Bourbaki, 1969/70, exposé n\textsuperscript{o} 376, Lecture Notes 180, Springer-Verlag.
\item[{[3]}] G. Harder -- \emph{Über die Galoiskohomologie halbeinfacher Matrizengruppen II}, Math. Zeitschrift, 92 (1966), p. 396--415.
\item[{[4]}] M. Kneser -- \emph{Schwache Approximation in algebraischen Gruppen}, Coll. sur la théorie des groupes algébriques, p. 41--52, Bruxelles 1962, CBRM.
\item[{[5]}] K. Miyake -- \emph{On models of certain automorphic function fields}.
\item[{[6]}] D. Mumford -- \emph{Geometric invariant theory}, Ergebnisse 34, 1965, Springer-Verlag.
\item[{[7]}] D. Mumford -- \emph{Families of abelian varieties}, in Alg. groups and discontinuous subgroups, p. 347--351, Boulder 1965, AMS.
\item[{[8]}] D. Mumford -- \emph{A Note on Shimura's Paper ``Discontinuous groups and Abelian Varieties''}. Math. Ann. 181, 345--351 (1969).
\item[{[9]}] V. P. Platonov -- \emph{Le problème d'approximation forte et l'hypothèse de Kneser--Tits}, Izvestia Acad. Nauk CCCP, 33 (1969), p. 1211--1219 and 34 (1970), p. 775--777.
\item[{[10]}] G. Shimura -- \emph{On analytic families of polarized abelian varieties and automorphic functions}, Ann. of Math., 78 1 (1963), p. 149--192.
\item[{[11]}] G. Shimura -- \emph{On the field of definition for a field of automorphic functions I, II, III}, Ann. of Math., 80 1 (1964), p. 160--189; 81 1 (1965), p. 124--165 and 83 2 (1966), p. 377--385.
\item[{[12]}] G. Shimura -- \emph{Construction of class fields and zeta functions of algebraic curves}, Ann. of Math. 85 1 (1967), p. 58--159.
\item[{[13]}] G. Shimura -- \emph{Algebraic number fields and symplectic discontinuous groups}, Ann. of Math., 86 3 (1967), p. 503--592.
\item[{[14]}] G. Shimura -- \emph{On canonical models of arithmetic quotients of bounded symmetric domains}, Ann. of Math., 91 1 (1970), p. 144--222.
\item[{[15]}] G. Shimura -- \emph{Automorphic functions and number theory}, Lecture Notes in Math. 54, 1968, Springer-Verlag.
\item[{[16]}] G. Shimura and T. Taniyama -- \emph{Complex multiplication of abelian varieties and its applications to number theory}, Publ. Math. Soc. Jap., 6 (1961).
\end{description}

\end{document}
